% Template created by Karol Kozioł (www.karol-koziol.net) for ShareLaTeX

\documentclass[a4paper,9pt]{extarticle}
\usepackage[utf8]{inputenc}
\usepackage[T1]{fontenc}
\usepackage{graphicx}
\usepackage{xcolor}
\usepackage{tikz}

\usepackage{amsmath,amssymb,textcomp}
\everymath{\displaystyle}

\usepackage{times}
\renewcommand\familydefault{\sfdefault}
\usepackage{tgheros}
\usepackage[defaultmono,scale=0.85]{droidmono}

\usepackage{multicol}
\setlength{\columnseprule}{0pt}
\setlength{\columnsep}{20.0pt}


\usepackage{geometry}
\geometry{
a4paper,
total={210mm,297mm},
left=10mm,right=10mm,top=10mm,bottom=15mm}

\linespread{1.3}


% custom title
\makeatletter
\renewcommand*{\maketitle}{%
\noindent
\begin{minipage}{0.65\textwidth}
\begin{tikzpicture}
\node[rectangle,rounded corners=6pt,inner sep=10pt,fill=blue!50!black,text width= 0.95\textwidth] {\color{white}\Huge \@title};
\end{tikzpicture}
\end{minipage}
\hfill

\bigskip\bigskip
}%
\makeatother

% custom section
\usepackage[explicit]{titlesec}
\newcommand*\sectionlabel{}
\titleformat{\section}
  {\gdef\sectionlabel{}
   \normalfont\sffamily\Large\bfseries\scshape}
  {\gdef\sectionlabel{\thesection\ }}{0pt}
  {
\noindent
\begin{tikzpicture}
\node[rectangle,rounded corners=3pt,inner sep=4pt,fill=blue!50!black,text width= 0.95\columnwidth] {\color{white}\sectionlabel#1};
\end{tikzpicture}
  }
\titlespacing*{\section}{0pt}{15pt}{10pt}


% custom footer
\usepackage{fancyhdr}
\makeatletter
\pagestyle{fancy}
\fancyhead{}
\fancyfoot[C]{\footnotesize \textcopyright\ \@date\ \ \@author}
\renewcommand{\headrulewidth}{0pt}
\renewcommand{\footrulewidth}{0pt}
\makeatother


\title{Cortex cheat sheet for bacterial genomics}
\date{2015}



\begin{document}

\maketitle

\begin{multicols*}{2}


\section{Requirements}

Download Cortex from github (https://github.com/iqbal-lab/cortex) as follows:
\begin{verbatim}
git clone --recursive https://github.com/iqbal-lab/cortex.git
\end{verbatim}
You must have installed VCFtools (and have the entire directory, not just the binary executable), R and Stampy. R must be in your path. Add the following directories to PERL5LIB and PATH
\begin{verbatim}
export PERL5LIB= /path/cortex/scripts/analyse_variants/
    bioinf-perl/lib;/path/cortex/scripts/calling:$PERL5LIB
export PATH = /path/cortex/scripts/analyse_variants/
   needleman_wunsch-0.3.0
\end{verbatim}

Compile and install cortex for 1,2,N,N+1 colours where N=number of samples

\section{Preparation (once per species)}
Assuming you want to report results with coordinates on a reference genome:
\begin{itemize}
\item{Build a Stampy hash of the genome}
\end{itemize}
\begin{verbatim}
stampy.py -G stampy_refhash ref.fa
stampy.py -g stampy_refhash -H stampy_refhash
\end{verbatim}
\begin{itemize}
\item{Build a Cortex binary graph file of the reference genome at the kmer size you are going to use (start with k=31). Don't worry about the --max\_read\_len 10000 argument, for reference fasta this just means the file is parsed in chunks of size 10000bp.}
\end{itemize}
\begin{verbatim}
ls /path/to/ref.fasta > ref_fa_list
cortex_var_31_c1 --kmer_size 31  --mem_height 17 
   --mem_width 100  --se_list ref_fa_list
   --max_read_len 10000  --dump_binary ref.k31.ctx 
   --sample_id REF
\end{verbatim}
\begin{itemize}
\item{Make an INDEX file which maps each sample to a LIST of fastq (or bams)}
\end{itemize}


\section{Compare each sample against a reference independently}

{\bf Step1 }First we build per-sample graphs in parallel using commands such as this GNU parallels command. Within directory OUTDIR1, we will make a directory for each sample. A typical command would be (for 1700 samples):
\begin{verbatim}
parallel --gnu -j 20  perl scripts/calling/par.pl 
  --num {}  --index INDEX   --list_ref LISTREF 
  --refbindir DIR_WITH_CTX_BINS --index_dir indexes/ 
  --stampy_bin stampy-1.0.23/stampy.py 
  --stampy_hash  stampy_refhash  --bc yes  --pd no 
  --kmer 31 --mem_height 18  --mem_width 100 
  --qthresh 10  --cortex_dir /path/to/cortex/ 
  --out_dir OUTDIR1  --vcftools_dir /path/vcftools_0.1.9/
    ::: {1..1700}
\end{verbatim}
{\bf Step2} Combine all the per-sample VCFs to get one combined set of sites (SNPs, indels, structural variants).
\begin{verbatim}
perl scripts/analyse_variants/combine/combine_vcfs.pl 
  --list_vcfs list_all_raw_vcfs  --vcftools_dir vcftools_0.1.9 
  --outdir OUTDIR2  --prefix presidents  --refname REF_v1 
  --ref_fasta ref.fa --rootdir_for_sample_output  OUTDIR1
  --kmer 31 --mem_height 18 --mem_width 150
   --ref_binary  ref.k31.ctx
\end{verbatim}
{\bf Step 3} Finally, independently genotype each sample at all of these sites. Here using GNU parallels to spread across 20 cores of a server. Memory use now drops as we only use the graph of poymorphisms.
\begin{verbatim}
cat OUTDIR2/list_args_for_final_step | parallel --colsep '\t' 
     perl scripts/calling/genotype_1sample_against_sites.pl 
	 --config OUTDIR2/config.txt 
 	--invcf OUTDIR2/presidents.sites_vcf 
	--sample {1} --outdir {2} --genome_size 24000000 
 	--sample_graph {3} --mem_height 17  --mem_width 100
\end{verbatim}
This will give you one VCF file per sample in OUTDIR1/\{sample\_id\}/union\_calls/

\section{Segregating variants within our dataset (joint workflow)}
First build sample graphs as Step1 in the previous example. Then UNDETERMINED



\section{Pan-genome analysis}

To detect presence of a set of predefined genes (genes.fasta) among your samples
\begin{verbatim}
\end{verbatim}
To look at pan-genome graph of all samples and see which samples have which contigs, allowing you to stratify them by frequency or look for differentiating/segregating contigs.


\section{The End}

For further information:


\end{multicols*}

\end{document}
